\begin{solapartat}
\punts{0,3} $h\nu + (-W_e) = E_c$

\punts{0,1} $h\nu$: Energia del fotó

\punts{0,1} $W_e$: Treball d'extracció de l'electró al metall

\punts{0,1} $E_c$: Energia cinètica màxima dels fotoelectrons

\medskip
\punts{0,4} La condició de llindar correspon a $E_c = 0$, per tant $h\nu_L = W_e$.

Com $\nu = \dfrac{c}{\lambda} \rightarrow \dfrac{hc}{\lambda_L} = W_e \Rightarrow \lambda_L = \dfrac{hc}{W_e}$

Nota: S'han d'acceptar com a vàlides expressions equivalents amb altres símbols per a la freqüència (\textit{f}) o el treball d'extracció (\textit{$\Phi$}), o expressant la $E_c$ com $(1/2)mv^2$ o en funció del potencial de frenat ($eV_\textit{frenat}$), o fins hi tot expressions com \underline{``energia} \underline{del} \underline{fotó} \underline{-} \underline{treball} \underline{d'extracció} \underline{de} \underline{l'electró} \underline{al} \underline{metall} \underline{=} \underline{energia} \underline{cinètica} \underline{màxima} \underline{dels} \underline{fotoelectrons''} (en aquest cas també s'explicita el significat de cada terme. En qualsevol cas, hi ha d'haver tres termes i cadascun ha de representar una energia.
\end{solapartat}

\begin{solapartat}
\punts{0,2} Energia del fotó: $h\nu = 4,375\times 10^{-19}\ \mathrm{J} = 2,735\ \mathrm{eV}$

% DUBTE: a l'original hi diu «1,60×19^{-19} J» (errata per 1,60×10^{-19} J).
\punts{0,5} $E_c = h\nu + (-W_e) = 4,375\times 10^{-19} - 2,36\ \mathrm{eV}\times\dfrac{1,60\times 19^{-19}\ \mathrm{J}}{1\ \mathrm{eV}} = 6,00\times 10^{-20}\ \mathrm{J}$

o bé, de manera equivalent

$E_c = h\nu + (-W_e) = 2,735\ \mathrm{eV} - 2,36\ \mathrm{eV} = 0,375\ \mathrm{eV}$

\medskip
\punts{0,3} Per frenar-los cal aplicar una diferència de potencial que provoqui un augment de l'energia potencial elèctrica igual a aquesta energia cinètica, per tant:
\[ \mathit{\Delta} V = 0,375\ \mathrm{V} \]

Nota: No s'exigeix a l'exercici que indiquin el signe
\end{solapartat}
